%==============================================================================
% DOCUMENTCLASS
%==============================================================================
\documentclass[12pt,oneside]{book}

%==============================================================================
% METADATOS
%==============================================================================
\newcommand{\universidad}{Universidad de Buenos Aires (UBA)}
\newcommand{\facultad}{Facultad de Derecho}
\newcommand{\carrera}{Carrera de Abogacía}
\newcommand{\materia}{Derecho Constitucional Profundizado}
\newcommand{\titulo}{Unidad introductoria --- El concepto de Constitución, sus fuentes y su interpretación}
\newcommand{\autor}{Isaac Edgar Camacho Ocampo}
\newcommand{\anio}{2026}

%==============================================================================
% PAQUETES
%==============================================================================
\usepackage[utf8]{inputenc}
\usepackage[T1]{fontenc}
\usepackage[provide=*,spanish]{babel}

\usepackage[
    top=2.5cm,
    bottom=2.5cm,
    left=3cm,
    right=2.5cm,
    headheight=15pt
]{geometry}

\usepackage{amsmath}
\usepackage{amssymb}
\usepackage{mathtools}

\usepackage{graphicx}
\usepackage{float}
\usepackage{caption}
\usepackage{subcaption}

\usepackage{booktabs}
\usepackage{array}
\usepackage{longtable}
\usepackage{tabularx}

\usepackage{enumitem}

\usepackage{xcolor}
\usepackage[most]{tcolorbox}

\usepackage{fancyhdr}
\usepackage{ragged2e}

\usepackage[
    colorlinks=true,
    linkcolor=blue!50!black,
    citecolor=green!40!black,
    urlcolor=blue!60!black,
    pdftitle={Apuntes de Derecho Constitucional Profundizado},
    pdfauthor={Isaac Edgar Camacho Ocampo},
    pdfsubject={Apuntes universitarios},
    bookmarks=true
]{hyperref}

%==============================================================================
% CONFIGURACIÓN
%==============================================================================
\graphicspath{
    {./img/}
    {./graficos/}
    {./figuras/}
}

\setcounter{tocdepth}{2}
\renewcommand{\contentsname}{Índice}

\setlength{\parindent}{0pt}
\setlength{\parskip}{0.5em}

\setlist[itemize]{
    topsep=0.3em,
    itemsep=0.2em
}

\setlist[enumerate]{
    topsep=0.3em,
    itemsep=0.2em
}

\clubpenalty=10000
\widowpenalty=10000

%==============================================================================
% COMANDOS
%==============================================================================
\newcommand{\abs}[1]{\left|#1\right|}
\newcommand{\norm}[1]{\left\lVert#1\right\rVert}
\newcommand{\vect}[1]{\mathbf{#1}}
\newcommand{\deriv}[2]{\frac{d#1}{d#2}}
\newcommand{\pderiv}[2]{\frac{\partial#1}{\partial#2}}

%==============================================================================
% ENTORNOS / CAJAS
%==============================================================================
\newtcolorbox{definition}[1]{
    enhanced,
    breakable,
    title={Definición: #1},
    fonttitle=\bfseries,
    colback=blue!3,
    colframe=blue!50!black,
    boxrule=0.8pt,
    arc=2mm
}

\newtcolorbox{important}{
    enhanced,
    breakable,
    title={Importante},
    fonttitle=\bfseries,
    colback=yellow!8,
    colframe=orange!70!black,
    boxrule=0.8pt,
    arc=2mm
}

\newtcolorbox{example}[1]{
    enhanced,
    breakable,
    title={Ejemplo: #1},
    fonttitle=\bfseries,
    colback=green!4,
    colframe=green!40!black,
    boxrule=0.8pt,
    arc=2mm
}

\newtcolorbox{formula}[1]{
    enhanced,
    breakable,
    title={Fórmula / Regla: #1},
    fonttitle=\bfseries,
    colback=purple!4,
    colframe=purple!50!black,
    boxrule=0.8pt,
    arc=2mm
}

\newtcolorbox{exercise}[1]{
    enhanced,
    breakable,
    title={Ejercicio: #1},
    fonttitle=\bfseries,
    colback=gray!5,
    colframe=gray!60!black,
    boxrule=0.8pt,
    arc=2mm
}

\newtcolorbox{note}{
    enhanced,
    breakable,
    title={Nota},
    fonttitle=\bfseries,
    colback=white,
    colframe=gray!50,
    boxrule=0.6pt,
    arc=2mm
}

% Caja adicional para transcribir textos normativos citados en el apunte
\newtcolorbox{norma}[1]{
    enhanced,
    breakable,
    title={#1},
    fonttitle=\bfseries,
    colback=white,
    colframe=black!65,
    boxrule=0.6pt,
    arc=2mm
}

%==============================================================================
% CONFIGURACIÓN FINAL
%==============================================================================
\pagestyle{fancy}
\fancyhf{}

\fancyhead[L]{\small\nouppercase{\leftmark}}
\fancyhead[R]{\small\materia}

\fancyfoot[C]{\thepage}

\renewcommand{\headrulewidth}{0.4pt}
\renewcommand{\footrulewidth}{0pt}

\fancypagestyle{plain}{
    \fancyhf{}
    \fancyfoot[C]{\thepage}
    \renewcommand{\headrulewidth}{0pt}
}

%==============================================================================
% DOCUMENT
%==============================================================================
\begin{document}

%------------------------------------------------------------------------------
% PORTADA
%------------------------------------------------------------------------------
\begin{titlepage}

\thispagestyle{empty}

\begin{center}

\vspace*{1cm}

{\large \textsc{\universidad}}

\vspace{0.2cm}

{\large \textsc{\facultad}}

\vspace{0.2cm}

{\large \carrera}

\vspace{3cm}

{\Huge \textbf{\materia}}

\vspace{1cm}

{\LARGE \titulo}

\vspace{0.8cm}

\textit{Comprender es el primer paso para convertir el estudio en conocimiento.}

\vspace{1.2cm}

\rule{0.70\textwidth}{0.5pt}

\vspace{0.5cm}

{\large Apuntes de estudio}

\vspace{0.5cm}

\rule{0.70\textwidth}{0.5pt}

\vfill

{\large \autor}

\vspace{0.3cm}

{\large \anio}

\end{center}

\vfill

\begin{flushleft}
\textit{Este es un modesto aporte a los estudiantes de ``Derecho Constitucional Profundizado'' no pretende sustituir el material oficial de la materia.}
\end{flushleft}

\end{titlepage}

%------------------------------------------------------------------------------
% ÍNDICE
%------------------------------------------------------------------------------
\tableofcontents

%------------------------------------------------------------------------------
% CONTENIDO
%------------------------------------------------------------------------------
\chapter*{Presentación}
\addcontentsline{toc}{chapter}{Presentación}
\markboth{Presentación}{Presentación}

Los siete bloques de estos apuntes funcionan como los cimientos, las paredes y el techo de un mismo edificio: la Constitución. Primero se presentan las reglas de convivencia del curso (el ``reglamento del edificio''); luego se define qué es una Constitución y de dónde proviene; después se analiza cómo se refuerza esa estructura con reformas y tratados sin derribarla; y finalmente se estudia cómo conviven dentro de ese edificio la Nación, las provincias y los municipios, y qué herramientas usan los intérpretes ---jueces, abogados, legisladores--- para decidir qué significa hoy un texto redactado hace más de 170 años.

\section*{¿Para qué sirve esta materia en la vida profesional?}

\textbf{¿Por qué importa más allá del examen?} Porque entrena una herramienta que se usa toda la vida profesional: la capacidad de argumentar con solidez frente a normas que no tienen una única respuesta correcta. El abogado litigante, el legislador, el negociador y el administrador enfrentan a diario ``zonas grises'' donde la ley no resuelve todo por sí sola, y ahí es donde se nota quién sabe razonar jurídicamente y quién no.

\textbf{¿Es útil para la vida, y no solo para la carrera?} Sí. Comprender cómo se interpreta la norma que organiza la convivencia de un país ---quién tiene poder, qué límites tiene ese poder y qué garantías protegen a cada persona--- permite participar como ciudadano informado en debates públicos (reformas, coparticipación, autonomía municipal, derechos) que afectan la vida cotidiana de cualquier habitante, sea o no abogado.

\textbf{¿Qué se lleva un futuro profesional de esta clase en particular?} Un mapa de las fuentes del derecho constitucional argentino, el manejo de los tres grandes métodos de interpretación (textualismo, originalismo y pragmatismo) y una primera aproximación al análisis económico del derecho: herramientas que se usan tanto para litigar como para asesorar, legislar o negociar.

%==============================================================================
\chapter[Encuadre de la materia]{Encuadre de la materia}
%==============================================================================

\section{Modalidad de cursada y evaluación}

La materia se dicta en modalidad virtual, pero los exámenes parciales y el examen final son \textbf{presenciales}. El docente manifestó su intención de compensar la distancia propia de la virtualidad con una mayor participación activa del estudiantado: propuso organizar exposiciones periódicas de los propios estudiantes, aproximadamente cada tres clases, en las que estos asuman un rol protagónico presentando temas trabajados de antemano. Para las futuras exposiciones y debates se conformarían equipos de entre tres y cuatro integrantes.

La cátedra prevé poner a disposición el programa de la materia y una bibliografía de base, complementada de forma progresiva con material adicional (fallos y artículos breves) a lo largo del cuatrimestre.

\subsection{Formato de evaluación}

La cursada contempla \textbf{dos exámenes parciales presenciales de carácter escrito}, con una modalidad mixta: preguntas de desarrollo más amplias combinadas con preguntas de respuesta más breve o de opción múltiple. El objetivo declarado de esta combinación es evaluar de forma más justa a estudiantes con distintos estilos de examen, evitando encasillar a nadie en un único sistema que lo favorezca o lo desfavorezca totalmente.

Los \textbf{recuperatorios}, en cambio, tendrían una modalidad más flexible ---trabajos, exposiciones u otras formas--- a definir según la cantidad de estudiantes que deban rendirlos y el avance general del curso.

\section{Objetivos y metodología de la cátedra}

La materia se encuentra en una instancia más avanzada de la carrera que el curso inicial de Derecho Constitucional, lo que permite dar por conocidos los grandes principios constitucionales para concentrarse en las denominadas \textbf{``zonas grises''}: los espacios de indeterminación normativa donde, a criterio del docente, se juega buena parte del desempeño profesional.

El objetivo central del curso fue explicado en estos términos: \textit{``nosotros en definitiva nos desenvolvemos más allá de ser abogado litigante, abogado legislador, abogado negociador, abogado administrador (...) siempre hay mucho de este constante choque de ideas y capacidad de argumentación''}. A partir de esa premisa, la finalidad pedagógica de la materia es dotar al estudiante de \textbf{pensamiento lateral} frente a los conflictos propios del derecho constitucional, junto con \textbf{capacidad crítica} y \textbf{capacidad de argumentación y persuasión}, entendidas como herramientas transferibles a cualquier campo profesional futuro.

\begin{important}
\textbf{El compromiso de honestidad intelectual.} La única regla de juego propuesta es la honestidad intelectual: cada estudiante puede expresar libremente su postura sobre los temas debatidos, ya que lo que se evalúa no es la coincidencia con una opinión determinada, sino la capacidad de articulación, la capacidad de argumentación y el manejo de conocimientos objetivos. La pluralidad de posturas enriquece la materia, en tanto disciplina social atravesada por la confrontación de ideas.
\end{important}

La participación activa ---debates, exposiciones grupales, confrontación de posturas entre dos grupos o formatos de ``órgano juzgador''--- es un componente central y recurrente de la cursada.

%==============================================================================
\chapter[El concepto de Constitución]{El concepto de Constitución}
%==============================================================================

\section{La Constitución como pacto social}

Como punto de partida se planteó una pregunta disparadora: qué se entiende por Constitución Nacional y qué cualidades debería tener.

Desde la ciencia política se aportó una mirada centrada en el diseño institucional, que ofrece una base valiosa para una aproximación más práctica al derecho constitucional, dado que la Constitución constituye el ``corpus político, económico, institucional y legal'' del país. Desde otra intervención se planteó una inquietud recurrente en la materia: por qué existe tanto temor a una reforma constitucional, cuando otros países modifican o amplían su texto de forma más continua; y se destacó la minuciosidad con que los constitucionalistas analizan ``la letra, la coma, el punto'' del texto constitucional.

Frente a la pregunta sobre qué es la Constitución, no existe una respuesta unívoca, sino distintas aproximaciones. Una de ellas es la caracterización propia del pensamiento de \textbf{Juan Bautista Alberdi}, para quien la Constitución es un verdadero \textbf{``plan político''}. La visión considerada predominante en la actualidad es la contractualista.

\begin{definition}{Constitución según el criterio contractualista}
Esquema institucional básico; acuerdo esencial que permite la vida en sociedad a partir de las restricciones que cada individuo acepta voluntariamente sobre parte de sus derechos, en pos de una armonización de la vida común dentro de un territorio, una soberanía y una población determinados.
\end{definition}

Esta idea se vincula con el \textbf{principio de gobierno representativo}: \textit{``cuando nosotros hablamos de que el pueblo no delibera ni gobierna, sino a través de sus representantes, también implica una limitación para el ejercicio de esa soberanía tan directa, tan propia (...) de la Revolución Francesa''}. Es decir, la Constitución no solo organiza al Estado, sino que también \textbf{autolimita el ejercicio directo de la soberanía popular}, canalizándola a través de instituciones representativas.

\begin{note}
\textbf{Precisión conceptual.} Se refirió el principio de gobierno representativo como propio de ``fines del siglo XVI, con la Revolución Francesa''. Corresponde precisar, por razones de exactitud histórica, que la Revolución Francesa tuvo lugar a fines del siglo XVIII (1789) y no del siglo XVI; se trata de una imprecisión menor, probablemente de expresión oral, que no afecta el argumento sustancial sobre la representación política.
\end{note}

\section{Dimensión económica de la Constitución}

La Constitución no se agota en su faz jurídico-institucional: también expresa un \textbf{programa económico} y un modelo de relación entre la Argentina y el resto del mundo ---con otros países, con los migrantes, con el capital humano y material que ingresa al país.

\textbf{Juan Bautista Alberdi} es una referencia ineludible del pensamiento económico incorporado a la Constitución, aun cuando fue más relevante ``desde el mundo de las ideas'' que desde el ejercicio de cargos políticos formales. Su obra \textit{Sistema Económico y Rentístico} fue mencionada como fuente de los principios del Estado liberal del siglo XIX que, a criterio del docente, la Constitución recoge en varios de sus preceptos esenciales, delineando así un verdadero diseño económico de la Nación en cuanto a las competencias del gobierno federal y su relación con las provincias.

Como cierre conceptual de este bloque se formuló una pregunta abierta sobre la titularidad de los recursos naturales, que funcionó como puente hacia el bloque federal desarrollado más adelante.

%==============================================================================
\chapter[Fuentes del constitucionalismo argentino]{Fuentes del constitucionalismo argentino}
%==============================================================================

\section{Influencias hispánica y americana}

Al no conservarse actas suficientemente concisas de la Convención Constituyente originaria, se han generado numerosas especulaciones acerca de cuál fue la base más relevante del texto de 1853.

La posición del docente fue clara: \textit{``mi respuesta es bastante clara y creo que sería lo más razonable hablar de una mixtura, pero por supuesto reconociendo la incidencia original de la Constitución americana en muchos puntos importantes''}. Junto con esa influencia estadounidense, también existieron aportes del \textbf{derecho hispano}, incluso cuando la intención de los constituyentes era abandonar esa impronta e inclinarse hacia un derecho considerado más ``original'' e ``incipiente'': el derecho americano. Ello ocurrió sobre todo por influencia de \textbf{Domingo Faustino Sarmiento}, aunque ---se remarcó--- también en cabeza de \textbf{Alberdi}, con su sólido dominio del derecho continental europeo y de ciertas herencias hispánicas.

\subsection{El corpus normativo de la materia}

El objeto de estudio de la asignatura es el conjunto de artículos que componen la Constitución Nacional, con la particularidad expansiva de los tratados internacionales incorporados por el \textbf{artículo 75, inciso 22}, calificado como ``famoso'' por su centralidad en el programa de la materia.

\section{El modelo constitucional de Estados Unidos}

Retomando la inquietud sobre el temor a las reformas constitucionales, se desarrolló el contraste entre el mecanismo de reforma estadounidense y el argentino.

Estados Unidos, formalmente, \textbf{nunca volvió a convocar una convención constituyente} para modificar su Constitución; sin embargo, muchos de sus principios más elementales ---como el \textbf{debido proceso}--- fueron incorporados mediante sucesivas \textbf{enmiendas}. Como ejemplo se citó la enmienda que estableció la prohibición de reelección presidencial indefinida, sancionada luego de los sucesivos períodos de Franklin D. Roosevelt: \textit{``se estableció la regla de la prohibición expresa; eso definió lógicamente el panorama político del país a posteriori''}.

En este marco se señaló que la garantía de no autoincriminación, popularmente conocida como ``la quinta enmienda'', tiene su equivalente en el artículo 18 de nuestra Constitución. Se confirmó que la mayoría de las grandes garantías del debido proceso parten de las primeras enmiendas estadounidenses.

\begin{norma}{Artículo 18 de la Constitución Nacional (garantía de no autoincriminación)}
\textit{``Ningún habitante de la Nación puede ser penado sin juicio previo fundado en ley anterior al hecho del proceso, ni juzgado por comisiones especiales, o sacado de los jueces designados por la ley antes del hecho de la causa. Nadie puede ser obligado a declarar contra sí mismo; ni arrestado sino en virtud de orden escrita de autoridad competente. Es inviolable la defensa en juicio de la persona y de los derechos.''}

\medskip
El artículo 18 es la norma que consagra en nuestro ordenamiento la garantía equivalente a la Quinta Enmienda estadounidense. El texto reproducido corresponde a la redacción vigente del artículo.
\end{norma}

\subsection{Trayectorias históricas inversas}

El contraste dio lugar a una reflexión histórica comparada entre ambos procesos constitucionales, que fueron trayectorias inversas: \textbf{Estados Unidos sancionó primero su Constitución} y luego atravesó una guerra civil de gran magnitud a mediados del siglo XIX; la \textbf{Argentina}, en cambio, transitó primero un largo período de guerras civiles ---de aproximadamente 30 a 35 años--- y recién después, con los estados provinciales y el estado nacional exhaustos, arribó a la sanción y consolidación de su Constitución, en particular a partir de 1860, con la incorporación de Buenos Aires.

\begin{table}[H]
\centering
\small
\renewcommand{\arraystretch}{1.25}
\begin{tabularx}{\textwidth}{
    >{\raggedright\arraybackslash}p{0.19\textwidth}
    >{\raggedright\arraybackslash}X
    >{\raggedright\arraybackslash}X
}
\toprule
\textbf{Aspecto} & \textbf{Estados Unidos} & \textbf{Argentina} \\
\midrule
Secuencia histórica &
Unión y Constitución primero; guerra civil después (mediados del siglo XIX) &
Guerras civiles primero (c.\ 30--35 años); consolidación constitucional después (1853--1860) \\

Mecanismo de reforma &
Enmiendas sucesivas aprobadas por el Congreso y ratificadas por las legislaturas estaduales &
Convención Constituyente formal; última reforma en 1994 \\

Grado de autonomía provincial/estadual originario &
Mayor autonomía relativa de los estados en el diseño federal inicial &
Tensión histórica entre gobierno federal fuerte y provincias, según el signo político de turno \\

Hito de incorporación territorial clave &
--- &
Incorporación de Buenos Aires a la Confederación (1860) \\
\bottomrule
\end{tabularx}
\caption{Comparación entre los procesos constitucionales de Estados Unidos y Argentina.}
\end{table}

\subsection{El Pacto de Olivos y la tendencia a romantizar los procesos constitucionales}

Se advirtió sobre la tendencia a \textbf{romantizar} los procesos constitucionales pasados. Así como hoy se identifica con facilidad la negociación política detrás del \textbf{Pacto de Olivos} ---antecedente inmediato de la reforma de 1994---, en 1853 los constituyentes también protagonizaron un proceso profundamente político. Como ejemplo se citó la incorporación de Buenos Aires en 1860, que respondió a un acuerdo político complejo: la provincia condicionó su ingreso a la Confederación a la modificación de determinados puntos que consideraba lesivos para sus intereses, en un contexto todavía marcado por las heridas recientes de la segunda batalla de Cepeda y la batalla de Pavón. En palabras del docente, ``no era simplemente una agitada discusión de café; había muertos todavía muy recientes''.

%==============================================================================
\chapter[Rigidez, reforma y tratados]{Rigidez, reforma y tratados internacionales}
%==============================================================================

\section{La reforma constitucional de 1994}

Los sistemas constitucionales suelen combinar un \textbf{núcleo duro de difícil modificación} ---la Constitución en sentido estricto--- con \textbf{procedimientos agravados} que buscan evitar reformas constantes, pero que a la vez admiten complementaciones relevantes con el correr del tiempo, tal como ocurrió en la práctica constitucional estadounidense mediante enmiendas.

\begin{important}
\textbf{El sistema argentino tras 1994.} La Constitución argentina, tras la reforma de 1994, presenta un mecanismo de incorporación de tratados internacionales que la enriquece de un modo más flexible que el rígido sistema de reforma total, aunque de ningún modo comparable en volumen a las enmiendas estadounidenses. Estas incorporaciones son ``solo complementariedades del cuerpo constitucional'' y nunca podrían contravenir la primera parte de la Constitución ---declaraciones, derechos y garantías--- ni su parte institucional.
\end{important}

\section{Tratados con jerarquía constitucional}

El mecanismo del \textbf{artículo 75, inciso 22} es una de las claves centrales del programa de la materia.

\begin{norma}{Artículo 75, inciso 22 de la Constitución Nacional}
\textit{``Corresponde al Congreso: (...) Aprobar o desechar tratados concluidos con las demás naciones y con las organizaciones internacionales y los concordatos con la Santa Sede. Los tratados y concordatos tienen jerarquía superior a las leyes. La Declaración Americana de los Derechos y Deberes del Hombre; la Declaración Universal de Derechos Humanos; la Convención Americana sobre Derechos Humanos; el Pacto Internacional de Derechos Económicos, Sociales y Culturales; el Pacto Internacional de Derechos Civiles y Políticos y su Protocolo Facultativo; la Convención sobre la Prevención y la Sanción del Delito de Genocidio; la Convención Internacional sobre la Eliminación de todas las Formas de Discriminación Racial; la Convención sobre la Eliminación de todas las Formas de Discriminación contra la Mujer; la Convención contra la Tortura y otros Tratos o Penas Crueles, Inhumanos o Degradantes; la Convención sobre los Derechos del Niño; en las condiciones de su vigencia, tienen jerarquía constitucional, no derogan artículo alguno de la primera parte de esta Constitución y deben entenderse complementarios de los derechos y garantías por ella reconocidos. Solo podrán ser denunciados, en su caso, por el Poder Ejecutivo Nacional, previa aprobación de las dos terceras partes de la totalidad de los miembros de cada Cámara. Los demás tratados y convenciones sobre derechos humanos, luego de ser aprobados por el Congreso, requerirán del voto de las dos terceras partes de la totalidad de los miembros de cada Cámara para gozar de la jerarquía constitucional.''}

\medskip
El texto reproducido corresponde a la redacción vigente del artículo 75, inciso 22, incorporado en la reforma constitucional de 1994.
\end{norma}

Hoy no existe discusión alguna respecto del \textbf{estatus constitucional} de estos tratados, ni de su aplicación directa y su exigibilidad: \textit{``se ha generado un mecanismo un poco más flexible por el cual la Constitución se ha visto aumentada, se ha visto robustecida con la incorporación más flexible del mecanismo de estos tratados''}. Estos tratados deben ser votados por una \textbf{mayoría agravada}, y la propia Constitución consagró de forma expresa la \textbf{prelación de los tratados internacionales por sobre las leyes del Congreso}, lo cual ya era sostenido previamente a nivel doctrinario.

Gracias a este mecanismo, el sistema constitucional argentino posterior a 1994 dejó de ser ``tan rígido tradicional como el que teníamos antes'', aunque ello no equivale al dinamismo del sistema de enmiendas estadounidense.

%==============================================================================
\chapter[Federalismo y recursos naturales]{Federalismo y recursos naturales}
%==============================================================================

\section{Dominio originario de los recursos naturales}

Retomando la dimensión económica de la Constitución, la titularidad de los recursos naturales permite ilustrar la relación entre el gobierno federal y las provincias. En términos básicos, la Constitución establece que los recursos naturales son de dominio público y, más concretamente, del \textbf{dominio originario de las provincias}.

\begin{norma}{Artículo 124 de la Constitución Nacional (in fine)}
\textit{``Corresponde a las provincias el dominio originario de los recursos naturales existentes en su territorio.''}

\medskip
La cláusula fue incorporada por la reforma constitucional de 1994 y consolidó a nivel constitucional el concepto de dominio originario provincial sobre los recursos naturales, previamente presente de modo implícito en normativa de rango inferior.
\end{norma}

\begin{example}{Vaca Muerta, ¿de quién es la riqueza del subsuelo?}
La formación de Vaca Muerta, situada entre las provincias de Neuquén y Mendoza, ilustra la disposición constitucional. La pregunta fue planteada en términos directos: ``¿la riqueza que hay en el sustrato de la provincia de Neuquén es de la provincia de Neuquén? Vaca Muerta es de Neuquén y Mendoza, o es de la Nación Argentina?''. Conforme al artículo 124, ese dominio originario corresponde a las provincias, lo cual genera un diseño institucional con fuerte impronta provincialista, incluso cuando la explotación y regulación efectiva de los recursos naturales en la práctica involucra también al gobierno federal.
\end{example}

\section{El régimen de coparticipación federal}

Se vinculó el punto anterior con la \textbf{ley de coparticipación federal}, norma ``recontravencida'' porque su plazo constitucional de sanción ya se encuentra vencido.

Para conjugar la idea de un dominio natural de las provincias que, a su vez, es regulado y en gran parte administrado en sus utilidades por el gobierno nacional, hay que pasar inevitablemente por esta ley: \textit{``lógicamente para poder conjugar la idea de un dominio natural de las provincias, pero que a su vez es regulado y es en gran parte administrada sus utilidades por el gobierno nacional, hay que pasar inevitablemente por el tamiz de esta bendita ley de coparticipación (...) que ya es una especie de Frankenstein que muestra sus cicatrices con el paso de los años''}. Su modificación o reemplazo es tan compleja, desde el punto de vista del consenso político necesario, que en la práctica el régimen se mantiene inalterado.

\begin{norma}{Disposición Transitoria Sexta de la Constitución Nacional}
\textit{``Un régimen de coparticipación conforme a lo dispuesto en el inciso 2 del artículo 75 y la reglamentación del organismo fiscal federal, serán establecidos antes de la finalización del año 1996; la distribución de competencias, servicios y funciones vigentes a la sanción de esta reforma, no podrá modificarse sin la aprobación de la provincia interesada; tampoco podrá modificarse en desmedro de las provincias la distribución de recursos vigente a la sanción de esta reforma y en ambos casos hasta el dictado del mencionado régimen de coparticipación.''}

\medskip
La reforma de 1994 impuso el plazo del 31 de diciembre de 1996 para sancionar un nuevo régimen de coparticipación conforme al artículo 75, inciso 2. A la fecha de la clase, dicho mandato constitucional continúa incumplido, lo que confirma la caracterización de la ley vigente como una norma ``recontravencida''.
\end{norma}

La discusión se amplió incorporando un actor adicional de relevancia creciente desde la reforma de 1994: la \textbf{constitucionalización de la autonomía municipal}, en particular en materia de estudios de impacto ambiental, donde el municipio interviene como actor político y jurídico relevante. Esta observación funcionó como puente hacia el bloque siguiente.

%==============================================================================
\chapter[Autonomía municipal]{Autonomía municipal}
%==============================================================================

\section{De la autarquía a la autonomía municipal}

A fines de los años noventa y principios de los 2000, el reconocimiento constitucional de los municipios era todavía una novedad reciente e incipiente, apenas consolidada por la reforma de 1994.

Un precedente jurisprudencial fundamental es el fallo \textbf{``Rivademar, Ángela Digna Balbina Martínez Galván de c. Municipalidad de Rosario''}, resuelto por la Corte Suprema de Justicia de la Nación el 21 de marzo de 1989 (Fallos 312:326). En dicho precedente, la Corte modificó su doctrina histórica y reconoció a los municipios el carácter de \textbf{entidades autónomas} ---y no meramente autárquicas---, sentando así una base jurisprudencial que la reforma constitucional de 1994 recogería expresamente en el artículo 123.

\begin{note}
En las notas originales, la mención de este precedente figuraba con una denominación errónea, producto de un error de reconocimiento de voz de la transcripción automática. Se consigna aquí la denominación corregida, manteniendo fiel el resto del relato.
\end{note}

\begin{norma}{Artículo 123 de la Constitución Nacional}
\textit{``Cada provincia dicta su propia constitución, conforme a lo dispuesto por el artículo 5° asegurando la autonomía municipal y reglando su alcance y contenido en el orden institucional, político, administrativo, económico y financiero.''}

\medskip
La cláusula de autonomía municipal fue incorporada por la reforma de 1994 y consolidó, a nivel constitucional, la doctrina sentada por la Corte Suprema en el fallo ``Rivademar'' (1989).
\end{norma}

Se describió, en términos de evolución institucional, el crecimiento sostenido del peso político y electoral de los municipios en las últimas décadas: \textit{``ha crecido mucho también el aspecto de la propia densidad de la burocracia municipal (...) ha crecido mucho la actividad y algunas competencias que están en cabeza de los municipios que antes no existían''}. Como ilustración de este fenómeno se mencionó la centralidad política actual de los intendentes del conurbano bonaerense en las elecciones de la provincia de Buenos Aires, un fenómeno que treinta años atrás prácticamente no existía como tal en el discurso público.

\section{Las tasas municipales y su desnaturalización}

Como contrapartida del crecimiento institucional de los municipios, se señaló el desarrollo paralelo de un sistema de financiamiento municipal cada vez más complejo y, en su caracterización, más agresivo hacia el contribuyente.

\begin{definition}{Tasa y tributo local}
Desde el ABC de la materia fiscal, la \textbf{tasa} debe guardar relación directa de \textbf{contraprestación} con un servicio efectivamente prestado por el municipio. Con el paso del tiempo esa noción se ha desvirtuado considerablemente: muchas tasas municipales se calculan de forma proporcional a determinadas actividades económicas o hechos imponibles ---terminología propia del derecho tributario--- sin relación directa con la contraprestación de un servicio concreto, asimilándose en la práctica a un verdadero \textbf{tributo local}.
\end{definition}

Este fenómeno se vincula con las quejas frecuentes sobre el carácter abusivo de ciertas tasas municipales y sobre la duplicación de tasas entre distintas jurisdicciones. Se concluyó: \textit{``el peso relativo de las tasas dentro de la tributación general ha crecido de forma considerable (...) es muy cierta tu afirmación en cuanto a ese nuevo actor político que son los municipios, que ha crecido mucho en cuanto a su impacto político, impacto electoral, impacto económico''}. El bloque se cierra con la relevancia creciente de los municipios en materia medioambiental como una de las competencias donde este nuevo protagonismo institucional se manifiesta con mayor claridad.

%==============================================================================
\chapter[Interpretación y análisis económico]{Interpretación constitucional y análisis económico}
%==============================================================================

\section{Métodos de interpretación constitucional}

Frente a las múltiples dimensiones de la Constitución ya expuestas, la materia se propone un abordaje \textbf{multidisciplinario} que no se agote en la terminología estrictamente jurídica. Se analizarán las llamadas ``consecuencias económicas de la interpretación constitucional'' como una herramienta más para enriquecer la capacidad de argumentación del futuro profesional.

El problema central de la interpretación constitucional es que un texto fundacional, a diferencia de las normas de rango inferior, tiende a permanecer en el tiempo sin replanteos automáticos, salvo que medie una reforma constitucional. En palabras del docente: \textit{``aquel pacto jurídico (...) la idea del contrato social, desde una visión ya sea rousseauniana o hobbesiana, a medida que van pasando los años se va quedando atrás, va quedando esa foto un poco más sepia de lo que era el inicio''}. Como ejemplo de esta aceleración de los cambios sociales frente a un texto estable, se mencionó que hace apenas diez años la inteligencia artificial se percibía como un tema de ciencia ficción y hoy plantea desafíos interpretativos concretos para el ordenamiento jurídico.

\subsection{Tres corrientes de interpretación constitucional}

Se presentaron tres grandes corrientes o escuelas de interpretación constitucional, ordenadas de mayor a menor apego literal al texto:

\begin{enumerate}
    \item \textbf{Textualismo}: aborda la interpretación como una disección gramatical del texto, buscando el significado corriente que las palabras tienen para la sociedad, evitando elucubraciones antojadizas.
    \item \textbf{Originalismo}: distinto del textualismo puro, se posiciona no tanto en la letra en sí misma sino en lo que los constituyentes quisieron legar o decir al sancionar la norma, lo que introduce un componente histórico-político con mayor margen de subjetividad.
    \item \textbf{Pragmatismo} (o interpretación dinámica): la corriente más flexible, que prioriza mantener incólumes y actualizados los principios constitucionales esenciales por sobre la voluntad explícita original del legislador, otorgando un rol más activo al intérprete.
\end{enumerate}

\begin{table}[H]
\centering
\small
\renewcommand{\arraystretch}{1.25}
\begin{tabularx}{\textwidth}{
    >{\raggedright\arraybackslash}p{0.17\textwidth}
    >{\raggedright\arraybackslash}X
    >{\raggedright\arraybackslash}X
}
\toprule
\textbf{Corriente} & \textbf{Fuente primaria de interpretación} & \textbf{Grado de subjetividad} \\
\midrule
Textualismo &
El texto constitucional y el significado corriente de sus palabras &
Relativamente bajo: busca una base objetiva en el sentido común del lenguaje \\

Originalismo &
La voluntad e intención histórica de los constituyentes &
Medio-alto: requiere reconstruir un contexto histórico-político \\

Pragmatismo &
Los principios constitucionales esenciales, actualizados a la época presente &
Alto: otorga mayor proactividad e iniciativa al intérprete \\
\bottomrule
\end{tabularx}
\caption{Corrientes de interpretación constitucional.}
\end{table}

El textualismo y el originalismo, aunque emparentados, \textbf{no son sinónimos}: \textit{``cabe decir que hay que diferenciarla del denominado originalismo (...) tiende a posicionarse más en lo que aquellos constituyentes nos quisieron legar, quisieron decir al sancionar esa Constitución''}, mientras que el textualismo se concentra en ``una especie de disección gramatical'' del texto mismo, sin indagar en la intención subjetiva de quienes lo redactaron.

\begin{important}
\textbf{La corriente pragmática y su mayor fertilidad interpretativa.} La posición pragmática es la que ofrece ``un campo mucho más fértil para reinterpretaciones constitucionales'', en tanto no busca su fundamento en la voluntad explícita del legislador original ni en el texto literal, sino en mantener vigentes y adaptados a los tiempos actuales los principios constitucionales esenciales, reconociendo a la vez que los constituyentes originales, aunque sabios, fueron ``producto de un tiempo que ya no existe''.
\end{important}

\section{Derecho y economía}

Como complemento de las herramientas de interpretación, la cátedra incorporará nociones de derecho y economía, y en particular de \textbf{análisis económico del derecho}, como recurso adicional para enriquecer la capacidad argumentativa de los estudiantes.

El alcance de este enfoque dentro de la materia fue delimitado con cuidado: \textit{``estamos en una materia que tiene sus límites y no somos especialistas en el tema, pero la idea de que detrás de todo gran diseño institucional hay también una posición que tiene sus efectos económicos concretos y sus dinámicas políticas que son muy concretas (...) eso sin duda no lo vamos a dejar pasar por alto''}. Cualquier análisis de un fallo, una decisión judicial o una decisión política debería tener razonablemente en cuenta estos postulados económicos.

\subsection{Cierre de la clase}

La idea que atravesó toda la clase es que la interpretación constitucional no es un ejercicio abstracto, sino una tarea cotidiana que se activa cada vez que se discute la constitucionalidad o inconstitucionalidad de una norma. Se anunció que en los encuentros siguientes se profundizaría en cada una de las corrientes interpretativas presentadas, así como en los fallos y debates de coyuntura que las ilustran.

%==============================================================================
% CORNELL
%==============================================================================
\chapter[Cuadro Cornell de repaso]{Cuadro Cornell de repaso}

La siguiente tabla organiza los conceptos clave de la clase bajo el formato Cornell: preguntas disparadoras a la izquierda y su desarrollo a la derecha, para facilitar el repaso activo del contenido.

\small
\renewcommand{\arraystretch}{1.3}
\begin{longtable}{
    >{\raggedright\arraybackslash}p{0.26\textwidth}
    >{\raggedright\arraybackslash}p{0.66\textwidth}
}
\toprule
\textbf{Preguntas / palabras clave} &
\textbf{Notas / respuestas} \\
\midrule
\endfirsthead

\toprule
\textbf{Preguntas / palabras clave} &
\textbf{Notas / respuestas} \\
\midrule
\endhead

\midrule
\multicolumn{2}{r}{\footnotesize\textit{Continúa en la página siguiente}} \\
\endfoot

\bottomrule
\endlastfoot

\textbf{¿Cuál es el objetivo central de la materia?} &
Desarrollar capacidad reflexiva, sana crítica, argumentación y persuasión frente a las ``zonas grises'' del derecho constitucional, más allá del rol específico que cada profesional termine ocupando (litigante, legislador, negociador, administrador). \\

\textbf{¿Cuál es la regla de juego propuesta para la cursada?} &
La honestidad intelectual: cada estudiante puede expresar libremente su postura, ya que lo que se evalúa no es la coincidencia con una opinión determinada, sino la capacidad de articulación, la capacidad de argumentación y el manejo de conocimientos objetivos. \\

\textbf{¿Cómo se evalúa la materia?} &
Con dos parciales presenciales escritos de modalidad mixta (preguntas de desarrollo y preguntas breves o de opción múltiple) y recuperatorios de formato más flexible según el caso. \\

\textbf{¿Qué es la Constitución según el criterio contractualista?} &
Un esquema institucional básico y un acuerdo esencial que permite la vida en sociedad a partir de restricciones voluntariamente aceptadas sobre los derechos individuales, con el fin de armonizar la convivencia en un territorio con soberanía y población propias. \\

\textbf{¿Cómo se vincula la Constitución con el gobierno representativo?} &
La Constitución no solo organiza al Estado, sino que también autolimita el ejercicio directo de la soberanía popular, canalizándola a través de instituciones representativas: el pueblo no delibera ni gobierna sino a través de sus representantes. \\

\textbf{¿Qué dimensión económica tiene la Constitución?} &
Expresa también un programa económico y un modelo de relación con el mundo (comercio, migración, capital), inspirado en buena medida en el pensamiento liberal del siglo XIX de Juan Bautista Alberdi y su obra \textit{Sistema Económico y Rentístico}. \\

\textbf{¿Cuáles son las principales fuentes históricas de la Constitución argentina?} &
Una mixtura con fuerte incidencia de la Constitución de los Estados Unidos, junto con aportes del derecho hispano e influencias del derecho continental europeo, sistematizadas especialmente por Alberdi y difundidas por Sarmiento. \\

\textbf{¿Qué diferencia el proceso constitucional argentino del estadounidense?} &
Estados Unidos sancionó su Constitución y luego atravesó su guerra civil; Argentina, en cambio, atravesó primero un largo período de guerras civiles y recién después consolidó su texto constitucional, con la incorporación de Buenos Aires en 1860 como hito clave. \\

\textbf{¿Qué garantía consagra el artículo 18 vinculada a la Quinta Enmienda estadounidense?} &
La garantía de que nadie puede ser obligado a declarar contra sí mismo, junto con el principio de juicio previo, juez natural, irretroactividad de la ley penal e inviolabilidad de la defensa en juicio. \\

\textbf{¿Qué establece el artículo 75, inciso 22?} &
Otorga jerarquía constitucional a un listado taxativo de tratados de derechos humanos ---ampliable por mayoría agravada de dos tercios de cada Cámara--- y consagra la prelación general de los tratados internacionales sobre las leyes del Congreso. \\

\textbf{¿Por qué el sistema constitucional argentino es menos rígido desde 1994?} &
Porque el mecanismo de incorporación de tratados con jerarquía constitucional permite robustecer el cuerpo normativo constitucional sin necesidad de una reforma total, aunque estas incorporaciones nunca pueden contravenir la primera parte de la Constitución. \\

\textbf{¿A quién pertenece el dominio originario de los recursos naturales?} &
A las provincias, conforme al artículo 124 de la Constitución Nacional, incorporado en la reforma de 1994 (ejemplo trabajado en clase: Vaca Muerta, entre Neuquén y Mendoza). \\

\textbf{¿Qué manda la Disposición Transitoria Sexta y se cumplió?} &
Ordenaba sancionar un nuevo régimen de coparticipación federal antes de finalizar 1996. El mandato continúa incumplido, lo que explica la caracterización de la ley vigente como una norma ``recontravencida''. \\

\textbf{¿Qué estableció el fallo ``Rivademar'' (CSJN, 21/03/1989)?} &
Modificó la doctrina histórica de la Corte y reconoció a los municipios como entidades autónomas ---y no meramente autárquicas---, sentando la base jurisprudencial que la reforma de 1994 recogería en el artículo 123. \\

\textbf{¿Qué diferencia conceptualmente a una tasa de un tributo local?} &
La tasa debe guardar relación directa de contraprestación con un servicio efectivamente prestado; cuando se calcula en proporción a una actividad económica sin esa contraprestación directa, se desnaturaliza y se asimila a un tributo local. \\

\textbf{¿Qué diferencia al textualismo del originalismo?} &
El textualismo interpreta el sentido gramatical corriente de las palabras del texto; el originalismo busca la intención histórica de los constituyentes al momento de sancionarlo, con mayor margen de subjetividad. \\

\textbf{¿En qué consiste la corriente pragmática de interpretación constitucional?} &
Prioriza mantener actualizados los principios constitucionales esenciales por sobre la letra o la intención histórica original, otorgando mayor proactividad al intérprete frente a un texto que envejece. \\

\textbf{¿Qué aporta el análisis económico del derecho a esta materia?} &
Una herramienta complementaria ---dentro de los límites de una materia jurídica, no económica--- para identificar las consecuencias y dinámicas económicas concretas que subyacen a todo gran diseño institucional o decisión judicial. \\

\midrule
\multicolumn{2}{p{0.94\textwidth}}{
\textbf{Resumen:}
La Constitución Nacional se presenta como un objeto de estudio con múltiples capas: es un pacto político que autolimita el ejercicio de la soberanía popular a través de la representación; un programa económico heredero del pensamiento liberal de Alberdi; un cuerpo normativo históricamente mixto, con influencias hispánicas y estadounidenses; y un diseño federal que reparte poder entre Nación, provincias y municipios, con tensiones en torno a los recursos naturales, la coparticipación y las tasas municipales. Desde 1994, un mecanismo más flexible de incorporación de tratados con jerarquía constitucional (artículo 75, inciso 22) robustece la Constitución sin necesidad de una reforma total. Frente a un texto que envejece sin actualizarse automáticamente, tres corrientes disputan su interpretación: el textualismo, apegado a la letra; el originalismo, atento a la voluntad histórica de los constituyentes; y el pragmatismo, orientado a mantener vigentes los principios esenciales. El análisis económico del derecho es una herramienta transversal para comprender las consecuencias reales de estas decisiones interpretativas. El curso se orienta a explorar las zonas grises donde la argumentación jurídica marca la diferencia.
} \\
\end{longtable}
\normalsize

\end{document}
